\section{Inference}
\label{sec:inference}

\subsection{Assumption}
\[\inferrule{ }{[P]}\]

\subsection{Conjunction}
\[\inferrule{P \\ Q}{P \wedge Q} \quad ({\wedge}I)\]
\[\inferrule{P \wedge Q}{P} \quad \inferrule{P \wedge Q}{Q} \quad ({\wedge}E)\]

\begin{align*}
    [P \wedge Q] \\
    P \\
    Q \\
    Q \wedge P
\end{align*}

\subsection{Implication}
\[\inferrule{P \\ P \Rightarrow Q}{Q} \quad ({\Rightarrow}E)\]
\[\inferrule{[P] \\\\ \vdots \\\\ Q}{P \Rightarrow Q}\]

\[(P \wedge Q) \Rightarrow (Q \wedge P)\]

\subsection{Classics}
\[\inferrule{\neg\neg P}{P} \quad ({\neg}{\neg} E)\]

\subsection{Quantifiers}
\[\inferrule{\forall a: P\langle a \rangle}{P\langle b \rangle} \quad ({\forall}E)\]
\[\inferrule{P\langle a \rangle}{\forall a: P\langle a \rangle} \quad ({\forall}I)\]
\[\inferrule{P\langle a \rangle}{\exists b \langle b \rangle} \quad ({\exists}I)\]
% Any notion of bound and free variables above is implicit - no further explanation of these concepts is provided.
